\documentclass[11pt,a4paper]{article}
% Shared preamble for RuPhO English problem statements (match T1 2026 style).
% Use: \input{latex/rupho_en_preamble.tex} after \documentclass.
\usepackage[margin=1in]{geometry}
\usepackage{amsmath,amssymb}
\usepackage{graphicx}
\usepackage{float}
\usepackage{enumitem}
\usepackage{microtype}
\usepackage{caption}
\usepackage{tabularx}
\usepackage{hyperref}

\captionsetup{font=small,labelfont=bf,skip=6pt}

\setlist[itemize]{leftmargin=1.2cm}
\setlist[enumerate]{leftmargin=1.2cm}

% One outer box per subproblem: [id | statement | points] (no inner rules)
\newcommand{\problembox}[3]{%
  \par\addvspace{0.42em}%
  \noindent
  \setlength{\fboxsep}{7.5pt}%
  \setlength{\fboxrule}{0.95pt}%
  \fbox{%
    \begin{minipage}{\dimexpr\linewidth-2\fboxsep-2\fboxrule\relax}
    \begin{tabularx}{\linewidth}{@{}>{\raggedright\arraybackslash\bfseries}p{2.1em} @{\hspace{0.9em}} >{\raggedright\arraybackslash}X @{\hspace{0.9em}} >{\raggedleft\arraybackslash\bfseries}p{2.4em}@{}}
    #1 & #2 & #3
    \end{tabularx}
    \end{minipage}%
  }%
  \par\addvspace{0.32em}%
}


\title{\textbf{Gaussian beams}}
\author{\normalsize Y25 (2025) --- English translation}
\date{}
\begin{document}
\maketitle

As is known, the propagation of electromagnetic waves is described by Maxwell's equations. Their simplest solutions in vacuum are plane waves. However, plane waves are not limited in space, and therefore are not physically realized. In this problem, beams that are symmetrical relative to the axis will be considered, in which the intensity decreases with distance to the beam axis $r$ according to the characteristic law $I \sim \exp(-2 r^2/w^2)$, where $w$ is the characteristic width of the beam. Such beams are called Gaussian beams. In such a beam, the field decreases quite quickly with distance from the axis, so that its width can be considered finite.  If the beam width is much greater than the wavelength, it is possible to obtain an analytical solution for the dependence of the wave field in the beam on coordinates.  It turns out that if in one of the sections of the beam the intensity dependence has a Gaussian form, then in all other sections it is also Gaussian. From Maxwell's equations it follows that each component of the electric and magnetic fields satisfies the wave equation, for example for $E_x$ it has the form \textbackslash\{\}[ \textbackslash\{\}frac\{\textbackslash\{\}partial^2 E_x\}\{\textbackslash\{\}partial x^2\}+\textbackslash\{\}frac\{\textbackslash\{\}partial^2 E_x\}\{\textbackslash\{\}partial y^2\}+\textbackslash\{\}frac\{\textbackslash\{\}partial^2 E_x\}\{\textbackslash\{\}partial z^2\} = \textbackslash\{\}frac\{1\}\{c^2\}  \textbackslash\{\}frac\{\textbackslash\{\}partial^2 E_x\}\{\textbackslash\{\}partial t^2\}. \textbackslash\{\}tag\{1\} \textbackslash\{\}] We will not consider other components in this problem. In the entire problem, we will consider only strictly monochromatic waves with a cyclic frequency $\omega$, which in vacuum corresponds to the wavelength $\lambda$ and wave number $k = 2\pi/\lambda$.  It is convenient to consider complex solutions of the wave equation of the form $$\textbackslash\{\}hat\{E\}_x (x, y,z,t)= v(x, y,z)\textbackslash\{\}exp(i \textbackslash\{\}omega t) ,$$ где $v$ – комплексная амплитуда волны в расcmатриваемой точке. Как и всегда, наблюдаемое поле связано с комплексным решением как: $$E_x(x, y, z, t) = \textbackslash\{\}mbox\{Re\} \textbackslash\{\}\{\textbackslash\{\}hat\{E_x\}(x, y, z, t)\textbackslash\{\}\}.$$ Если $\hat{E_x}$ удовлетворяет волновому уравнению, то в силу линейности  $E_x$ также удовлетворяет волновому уравнению.

Из уравнений Максвелла следует, что каждая компонента электрического и магнитного полей удовлетворяет волновому уравнению, например для $E_x$ оно имеет вид

\textbackslash\{\}[ \textbackslash\{\}frac\{\textbackslash\{\}partial^2 E_x\}\{\textbackslash\{\}partial x^2\}+\textbackslash\{\}frac\{\textbackslash\{\}partial^2 E_x\}\{\textbackslash\{\}partial y^2\}+\textbackslash\{\}frac\{\textbackslash\{\}partial^2 E_x\}\{\textbackslash\{\}partial z^2\} = \textbackslash\{\}frac\{1\}\{c^2\}  \textbackslash\{\}frac\{\textbackslash\{\}partial^2 E_x\}\{\textbackslash\{\}partial t^2\}. \textbackslash\{\}tag\{1\} \textbackslash\{\}]

Другие компоненты в этой задаче расcmатривать не будем. Во всей задаче мы будем расcmатривать только строго монохроматические волны с циклической частотой $\omega$, которой в вакууме отвечает длина волны $\lambda$ и волновое число $k = 2\pi/\lambda$.  Удобно расcmатривать комплексные решения волнового уравнения вида

$$\textbackslash\{\}hat\{E\}_x (x, y,z,t)= v(x, y,z)\textbackslash\{\}exp(i \textbackslash\{\}omega t) ,$$

где $v$ – комплексная амплитуда волны в расcmатриваемой точке. Как и всегда, наблюдаемое поле связано с комплексным решением как: $$E_x(x, y, z, t) = \textbackslash\{\}mbox\{Re\} \textbackslash\{\}\{\textbackslash\{\}hat\{E_x\}(x, y, z, t)\textbackslash\{\}\}.$$ Если $\hat{E_x}$ удовлетворяет волновому уравнению, то в силу линейности  $E_x$ also satisfies the wave equation.

\section*{Part A. Diffraction of Gaussian beams (2.0 points)}

In this part we will obtain various estimates for a Gaussian beam using elementary methods from diffraction theory. We assume that the beam propagates along the $z$ axis, the field oscillations in the $z = 0$ plane occur in phase, and the dependence of the amplitude on the distance $r = \sqrt{x^2 + y^2}$ to the symmetry axis of the system has the form $$ v(x, y, 0) = E_0 \textbackslash\{\}exp(- r^2/ w_0^2), $$где $E_0$ – вещественная постоянная. Будем искать распределение амплитуды волны по экрану, расположенному в плоскости $z = L$.

\problembox{A1}{At very large distances $L$, estimate the size of the spot $D$ on the screen by order of magnitude. Express your answer in terms of $\lambda$, $w_0$, $L$.}{0.30}

To calculate the diffraction pattern, we will use the Kirchhoff integral, which expresses the complex amplitude of the wave at any point in the half-space $z > 0$ in the form of an integral over the surface $z = 0$: $$ v(x, y, z) = - \textbackslash\{\}frac\{1\}\{i \textbackslash\{\}lambda\} \textbackslash\{\}int dx' dy'\textbackslash\{\},  \textbackslash\{\}frac\{v(x', y', 0) \}\{s\} e^\{-i k s\}, $$где $s$ – расстояние от точки плоскости до точки наблюдения.  Даже этот интеграл нельзя вычислить точно, поэтому вам потребуется использовать приближения, отвечающие случаям дифракции Фраунгофера и Френеля. Положение точки на экране будем характеризовать расстоянием $r = \sqrt{x^2 +y^2}$ до оси сиmmетрии системы, причем будем считать $r \ll L$. Let us first consider the case of Fraunhofer diffraction. In this case, the distance from the plane to the observation point is so great that all rays going to the observation point can be considered parallel.

\problembox{A2}{Under what constraint on $L$ is Fraunhofer diffraction theory applicable? Express your answer in terms of $L$, $w_0$, $\lambda$.}{0.20}

\problembox{A3}{In this case, write down the expression for the amplitude of the electric field in the form of an integral over $dx'dy'$. Find an expression for the field amplitude on the screen depending on the distance to the symmetry axis $r$. The answer may also include $E_0$, $w_0$, $k$. You may need the integral $$ \textbackslash\{\}int\textbackslash\{\}limits_\{-\textbackslash\{\}infty\}^\{+\textbackslash\{\}infty\}  dx\textbackslash\{\}, e^\{- ax^2 + i b x\} = \textbackslash\{\}sqrt\{\textbackslash\{\}frac\{\textbackslash\{\}pi\}\{a\}\}e^\{- b^2/4 a\}. $$}{0.80}

\problembox{A4}{Now consider the case of Fresnel diffraction. In this case, the distance $L$ is arbitrary, however, the angles between the rays going to the observation point and the axis $z$ can be considered small. Find $A(z) =  v(0, 0, z) \exp\left(ikz\right)$ – the complex amplitude on the $Oz$ axis without taking into account the phase of the plane wave. Express your answer in terms of $E_0$, $w_0$, $k$, $z$.}{0.70}

\section*{Part B. Gaussian beam as a solution to the wave equation (2.4 points)}

Since the Gaussian beam propagates along the $z$ axis, its amplitude contains a rapidly changing factor $e^{-ikz}$, we will look for a field in the form \textbackslash\{\}[ \textbackslash\{\}hat\{E\}_x(x, y, z, t) = u(x, y, z) \textbackslash\{\}exp(-i k z) \textbackslash\{\}exp(i \textbackslash\{\}omega t), \textbackslash\{\}tag\{2\}\textbackslash\{\}] with $\omega = ck$, and in the plane $z = 0$ $$ u(r, 0) =E_0 \textbackslash\{\}exp(- r^2/w_0^2). $$

\problembox{B1}{Substitute $(2)$ into $(1)$ and get the differential equation satisfied by $u(x, y, z)$.}{0.30}

Let us assume that $u(x, y, z)$ is a function that varies little along $Oz$ at a distance $\sim \lambda$. This means that $(2)$ describes a plane wave modulated by a function that varies slowly over distances on the order of the wavelength. Mathematically, this condition is written as: \textbackslash\{\}[ \textbackslash\{\}left|\textbackslash\{\}frac\{\textbackslash\{\}partial^2 u\}\{\textbackslash\{\}partial z^2\}\textbackslash\{\}right| \textbackslash\{\}ll \textbackslash\{\}left|k \textbackslash\{\}frac\{\textbackslash\{\}partial u\}\{\textbackslash\{\}partial z\}\textbackslash\{\}right|. \textbackslash\{\}tag\{\#\}\textbackslash\{\}] In this approximation, the equation of item B1 takes the form: \textbackslash\{\}[ \textbackslash\{\}frac\{\textbackslash\{\}partial^2 u\}\{\textbackslash\{\}partial x^2\} + \textbackslash\{\}frac\{\textbackslash\{\}partial^2 u\}\{\textbackslash\{\}partial y^2\} - 2 i k \textbackslash\{\}frac\{\textbackslash\{\}partial u\}\{\textbackslash\{\}partial z\} = 0. \textbackslash\{\}tag\{3\} \textbackslash\{\}]

As before $r = \sqrt{x^2 + y^2}$. The Gaussian beam corresponds to the solution of the equation $(3)$ of the form: \textbackslash\{\}[ u(r, z) = A(z) \textbackslash\{\}exp \textbackslash\{\}left(-\textbackslash\{\}frac\{i k r^2\}\{2 q(z)\} \textbackslash\{\}right), \textbackslash\{\}tag\{4\}\textbackslash\{\}]where $A(z)$, $q(z)$ are complex-valued functions, and the function $A(z)$ is the complex amplitude of the wave on axis of symmetry found in the previous part. Later in this part you will define the function $q(z)$.

\problembox{B2}{Show that the sum of the second partial derivatives $u$ with respect to the “transverse” coordinates $x$, $y$ can be represented as: $$\textbackslash\{\}frac\{\textbackslash\{\}partial^2 u\}\{\textbackslash\{\}partial x^2\} + \textbackslash\{\}frac\{\textbackslash\{\}partial^2 u\}\{\textbackslash\{\}partial y^2\} = \textbackslash\{\}left[-2 i \textbackslash\{\}chi(z) - \textbackslash\{\}chi^2(z) \textbackslash\{\}cdot r^2 \textbackslash\{\}right]\textbackslash\{\}cdot u.$$ Выразите $\chi(z)$ через $q(z)$, $k$.}{0.50}

\problembox{B3}{Show that the partial derivative of $u$ with respect to $z$ can be represented as: $$\textbackslash\{\}frac\{\textbackslash\{\}partial u\}\{\textbackslash\{\}partial z\} = \textbackslash\{\}left[\textbackslash\{\}frac\{A'(z)\}\{A(z)\}  + i \textbackslash\{\}zeta(z) \textbackslash\{\}cdot r^2 \textbackslash\{\}right] \textbackslash\{\}cdot u.$$Выразите $\zeta(z)$ через $q(z)$, $q'(z)$, $k$.}{0.30}

Substituting the results of items B1, B2 into $(3)$ and reducing to $u$ leads to the relation: \textbackslash\{\}[-2 i k \textbackslash\{\}frac\{A'(z)\}\{A(z)\} - 2 i \textbackslash\{\}chi(z) + \textbackslash\{\}left[2 k \textbackslash\{\}zeta(z) - \textbackslash\{\}chi^2(z) \textbackslash\{\}right] \textbackslash\{\}cdot r^2 = 0. \textbackslash\{\}tag\{5\}\textbackslash\{\}] A necessary condition for its implementation is the independence of the left side from $r$, that is, $$2k \textbackslash\{\}zeta(z) - \textbackslash\{\}chi^2(z) = 0.$$

\problembox{B4}{Derive a differential equation on $q(z)$. Solve it with the initial condition $$q(0) = i z_0,$$ где $z_0$ – действительная величина, $z_0 > 0$. Выразите $z_0$ через $w_0$, $k$.}{0.40}

Note that the approximation $(\#)$ will be true under the condition $k z_0 \gg 1$, which we will assume to be satisfied everywhere below. Let $$\textbackslash\{\}frac\{1\}\{q(z)\} = \textbackslash\{\}frac\{1\}\{R(z)\} - i\textbackslash\{\}frac\{2\}\{kw^2(z)\},$$где $R(z)$, $w(z)$ be real-valued functions.

\problembox{B5}{Express $R(z)$ and $w(z)$ in terms of $k$, $z$, $z_0$.}{0.20}

\problembox{B6}{Check that the function $A(z)$ found in Part A really satisfies the equation $(5)$, that is, when substituting it, the terms not containing $r^2$ become 0.}{0.40}

\problembox{B7}{Let us represent the function $A(z) $ as $$ A(z) = A_0 (z) e^\{i \textbackslash\{\}psi (z)\} $$с вещественными $A_0$ и $\psi$. Выразите $A_0$ и $\psi$ через $z$, $z_0$, $E_0$.}{0.30}

\section*{Part C. Characteristics of a Gaussian beam (4.0 points)}

In the approximation under consideration, the light intensity at an arbitrary point is equal to $$I = \textbackslash\{\}beta \textbackslash\{\}langle E^2 \textbackslash\{\}rangle,$$ где $\beta$ – константа, $\langle E^2 \rangle$ – средний по времени квадрат напряжённости поля в этой точке. Также в расcmатриваемом приближении энергия поля переносится волной вдоль оси $Oz$, there is no energy transfer in the perpendicular direction.

\problembox{C1}{Find the intensity distribution in the beam $I(r, z)$. Express your answer in terms of $\beta$, $A_0(z)$, $w(z)$. Draw a characteristic graph of the intensity distribution in the radial direction for a fixed $z$.}{0.50}

\problembox{C2}{Find the total radiation power $P$ transmitted by the beam through the plane $z = \operatorname{const}$. Express your answer in terms of $\beta$, $E_0$, $k$, $z_0$, $z$.}{0.40}

We will call the beam boundary the points at which the intensity is $e^2$ times less than at a point with the same coordinate $z$ and lying on the beam axis. More formally: if we consider a point with coordinate $z$ along the axis $Oz$ and removed at a distance $r$ from it, then it will lie on the boundary if: $$\textbackslash\{\}frac\{I(r, z)\}\{I(0, z)\} = \textbackslash\{\}frac\{1\}\{e^2\}.$$ В силу осевой сиmmетрии решения любая плоскость $z = \operatorname{const}$ пересекает границу по окружности. Диаметр этой окружности $d(z)$ we will call the width of the beam.

\problembox{C3}{Express the beam width $d(z)$ in terms of $w(z)$. Construct a qualitative graph of the dependence $d(z)$.}{0.40}

\problembox{C4}{Find the minimum beam width $d_п$. Express your answer in terms of $z_0$, $k$. The quantity $d_п$ is called the waist width. The plane perpendicular to $Oz$, in which the beam width is equal to the waist width, is called the waist plane.}{0.10}

A wave front is a surface with a constant phase of a wave. A single wave front passes through any point in space.

\problembox{C5}{Derive an equation that determines the shape of the wavefront of a Gaussian beam passing through a given point $(0, z_1)$. The equation may contain $R(z)$, $\psi(z)$, $k$, $z$, $z_1$.}{0.40}

As stated earlier, $kz_0 \gg 1$. Taking this into account, we can consider the functions $R(z)$ and $\psi(z)$ constant at small distances from the axis $Oz$. In this approximation, the shape of the wave front of a Gaussian beam coincides with the shape of the front of a spherical wave observed at a certain distance $\rho(z)$ from the source. Naturally, $\rho(z_1)$ is the radius of curvature of the wavefront of the Gaussian beam at point $(0,0,z_1)$.

\problembox{C6}{Find the radius of curvature of the front $\rho(z_1)$. Express your answer in terms of $R(z)$. Construct a graph of the dependence $\rho(z_1)$, indicating the characteristic features and their coordinates.}{0.50}

\problembox{C7}{Let at some point $z = z_a$ the parameters of the Gaussian beam $R(z_a) = R_a$, $w(z_a) = w_a$. Find $z_a$ and the beam waist diameter $d_п$. Express your answer in terms of $R_a$, $w_a$, $k$.}{0.70}

\problembox{C8}{Let a Gaussian beam with wavelength $\lambda_0 = 10.6~um$ have widths $d_1 = 3.398~mm$ and $d_2 = 6.760~mm$ at two points, the distance between which is $L = 10~cm$. Numerically calculate the width of the beam waist $d_п$.}{1.00}

\section*{Part D. Gaussian beams in lenses (1.6 points)}

Let us first consider the transformation of a plane wave by a thin lens. Let the plane of a thin converging lens with focal length $f$ be perpendicular to the axis $Oz$, the center of the lens is at point $(0, 0, z_л)$. A plane wave falls on the lens, one of the field components in which is: $$\textbackslash\{\}hat\{E\}_\{пx\}(x, y, z, t) = \textbackslash\{\}hat\{E\}_\{0x\} \textbackslash\{\}exp(i (\textbackslash\{\}omega t - k z)).$$ Известно, что после прохождения линзы сразу за ней напряженность поля оказывается равной: $$\textbackslash\{\}hat\{E'\}_\{пx\} (x, y, z_л, t) = \textbackslash\{\}hat\{E\}_\{пx\} (x, y, z_л, t) \textbackslash\{\}cdot \textbackslash\{\}hat\{\textbackslash\{\}tau\}(x, y),$$ здесь $\hat{\tau}(x,y)$ – коэффициент пропускания линзы. Известно, что для линзы он имеет вид $$\textbackslash\{\}hat\{\textbackslash\{\}tau\}(x, y) = \textbackslash\{\}exp(i \textbackslash\{\}eta (x^2+y^2)),$$ $\eta = \operatorname{const}$.

\problembox{D1}{It is known that a lens changes the phase of the incident wave in such a way that all rays come to focus with the same phase. Using this fact, express the coefficient $\eta$ in terms of $k$, $f$.}{0.70}

For Gaussian beams, the same transformation law can be approximately applied.

\problembox{D2}{Let a thin collecting lens with focal length $f$ be hit by a Gaussian beam propagating along $Oz$. The lens plane is perpendicular to the $Oz$ axis. Immediately in front of the lens, the beam parameters are $R(z_л) = R_1$, $w(z_л) = w_1$. It is known that after the lens it will also remain Gaussian.  Find the beam parameters $R_2$, $w_2$ immediately after passing through the lens. The lens diameter is much larger $w_1$.}{0.50}

\problembox{D3}{Find a condition connecting $R_1$, $w_1$ and $f$, under which the Gaussian beam will be focused behind the lens, that is, its waist plane will be located after the lens.}{0.40}

\vspace{1em}
\noindent\textit{Source:} \url{https://pho.rs/p/4222}

\end{document}